%!TEX TS-program = xelatex
%!TEX encoding = UTF-8 Unicode
% Awesome CV LaTeX Template for CV/Resume
%
% This template has been downloaded from:
% https://github.com/posquit0/Awesome-CV
%
% Author:
% Claud D. Park <posquit0.bj@gmail.com>
% http://www.posquit0.com
%
% Template license:
% CC BY-SA 4.0 (https://creativecommons.org/licenses/by-sa/4.0/)
%


%-------------------------------------------------------------------------------
% CONFIGURATIONS
%-------------------------------------------------------------------------------
% A4 paper size by default, use 'letterpaper' for US letter
\documentclass[11pt, a4paper]{awesome-cv}
\usepackage{tikz}
\usepackage{xcolor}
% Configure page margins with geometry
\geometry{left=1.4cm, top=.8cm, right=1.4cm, bottom=1.8cm, footskip=.5cm}

% Specify the location of the included fonts
\fontdir[fonts/]

% Color for highlights
% Awesome Colors: awesome-emerald, awesome-skyblue, awesome-red, awesome-pink, awesome-orange
%                 awesome-nephritis, awesome-concrete, awesome-darknight
\colorlet{awesome}{awesome-concrete}
% Uncomment if you would like to specify your own color
% \definecolor{awesome}{HTML}{CA63A8}

% Colors for text
% Uncomment if you would like to specify your own color
% \definecolor{darktext}{HTML}{414141}
% \definecolor{text}{HTML}{333333}
% \definecolor{graytext}{HTML}{5D5D5D}
% \definecolor{lighttext}{HTML}{999999}

% Set false if you don't want to highlight section with awesome color
\setbool{acvSectionColorHighlight}{false}

% If you would like to change the social information separator from a pipe (|) to something else
%\renewcommand{\acvHeaderSocialSep}{\quad\textbar\quad}
\renewcommand{\acvHeaderSocialSep}{\quad•\quad}


%-------------------------------------------------------------------------------
%	PERSONAL INFORMATION
%	Comment any of the lines below if they are not required
%-------------------------------------------------------------------------------
% Available options: circle|rectangle,edge/noedge,left/right
% \photo{./examples/profile.png}
\name{Filipe}{Guimarães}
\position{Msc Software Engineering Student}
\address{Braga, Portugal}

\mobile{(+351) 916037364}
\email{filipe\_guimaraes@live.com.pt}
%\homepage{filipeguimaraes.github.io}
\github{filipeguimaraes}
\linkedin{filipeguimaraes}
\gitlab{filipeguimaraes}
% \stackoverflow{SO-id}{SO-name}
% \twitter{@twit}
% \skype{skype-id}
% \reddit{reddit-id}
% \medium{madium-id}
% \googlescholar{googlescholar-id}{name-to-display}
%% \firstname and \lastname will be used
% \googlescholar{googlescholar-id}{}
% \extrainfo{extra informations}

%\quote{``Be the change that you want to see in the world."}


%-------------------------------------------------------------------------------
\begin{document}

% Print the header with above personal informations
% Give optional argument to change alignment(C: center, L: left, R: right)
%\makecvheader[L]
  \begin{minipage}[h]{.5\textwidth}
    %\raggedright
    \makecvheader[L]
  \end{minipage}% <---------------- Note the use of "%"
  \begin{minipage}[h]{.5\textwidth}
    %\raggedright
    %\vspace{1pt}
    %\includegraphics[width=0.5\textwidth]{photo.jpg}
    \hfill
    \begin{tikzpicture}
      \clip (0,0)  circle (2cm) ;
      \node[anchor=center] at (0,0) {\includegraphics[width=4cm]{photo.jpg}};
      %adjust this coordinate to move image
    \end{tikzpicture}

  \end{minipage}


% Print the footer with 3 arguments(<left>, <center>, <right>)
% Leave any of these blank if they are not needed
\makecvfooter
{\today}
{Filipe Guimarães~~~·~~~Curriculum Vitae}
{\thepage}


%-------------------------------------------------------------------------------
%	CV/RESUME CONTENT
%	Each section is imported separately, open each file in turn to modify content
%-------------------------------------------------------------------------------
%
\cvsection{About me}

HI.


%-------------------------------------------------------------------------------
%	SECTION TITLE
%-------------------------------------------------------------------------------
\cvsection{Education}


%-------------------------------------------------------------------------------
%	CONTENT
%-------------------------------------------------------------------------------
\begin{cventries}

  %---------------------------------------------------------
  \cventry
  {MSc Software Engineering (Mestrado integrado em Engenharia Informática)} % Degree
  {Universidade do Minho} % Institution
  {Braga, Portugal} % Location
  {Set. 2017 - now} % Date(s)
  {
    \begin{cvitems} % Description(s) bullet points
      \item { Bachelor in software Engineering. }
      \item { Master with focus on computer networks and network management. Some complementary classes on Data Mining and Computer Vision. }
      \item { Currently, developing my master thesis titled "\textit{QUIC Protocol Assessment in Adverse Network Conditions}".}
    \end{cvitems}
  }

  %---------------------------------------------------------
\end{cventries}
\begin{minipage}[t]{.53\textwidth}
%-------------------------------------------------------------------------------
%	SECTION TITLE
%-------------------------------------------------------------------------------
\cvsection{Professional Experience}


%-------------------------------------------------------------------------------
%	CONTENT
%-------------------------------------------------------------------------------
\begin{cventries}

%---------------------------------------------------------
\cventry
  {Trainee} % Job title
  {DSTelecom} % Organization
  {Braga, Portugal} % Location
  {Sept. 2022 - Dec. 2022} % Date(s)
  {
    \begin{cvitems} % Description(s) of tasks/responsibilities
      \item {Curricular internship with the aim to develop my thesis. Won some important insights of working in a company.}
    \end{cvitems}
}


%---------------------------------------------------------
\end{cventries}

%-------------------------------------------------------------------------------
%	SECTION TITLE
%-------------------------------------------------------------------------------
\cvsection{Other Experience}


%-------------------------------------------------------------------------------
%	CONTENT
%-------------------------------------------------------------------------------
\begin{cventries}
  \cventry
  {Colaborator} % Job title
  {\href{https://stolscience.com/}{STOL-Science Through Our Lives}} % Organization
  {Braga, Portugal} % Location
  {2023 - now} % Date(s)
  {
    \begin{cvitems} % Description(s) of tasks/responsibilities
      \item {Colaborator in the project "STOL-Science Through Our Lives", which aims to be a science communication group from different areas.}
      \item {Currently finishing a mobile aplicattion to preserve the historic information of the live stock observerd in the community garden of University of Minho.}
    \end{cvitems}
}

\cventry
  {Staff Member} % Job title
  {SEI \& JOIN} % Organization
  {Braga, Portugal} % Location
  {2019 - 2022} % Date(s)
  {
    \begin{cvitems} % Description(s) of tasks/responsibilities
      \item {STAFF Member for the "Semana de Engenharia Informárica" and "Jornada de informática".}
    \end{cvitems}
}

  \cventry
  {Staff Member} % Job title
  {ENEI} % Organization
  {Braga, Portugal} % Location
  {2020} % Date(s)
  {
    \begin{cvitems} % Description(s) of tasks/responsibilities
      \item {STAFF Member for the "Encontro nacional de Engenharia Informática".}
    \end{cvitems}
  }

\cventry
    {Coordinator} % Job title
    {Criar Asas} % Organization
    {Braga, Portugal} % Location
    {2018 - 2021} % Date(s)
    {
      \begin{cvitems} % Description(s) of tasks/responsibilities
        \item {Coordinator of the \textit{Criar Asas} youth group, where I organized activities to help the surrounding community.}
      \end{cvitems}
    }
%---------------------------------------------------------
\end{cventries}

\end{minipage}
\hfill
\begin{minipage}[t]{.45\textwidth}
%-------------------------------------------------------------------------------
%	SECTION TITLE
%-------------------------------------------------------------------------------
\cvsection{Academic Projects}


%-------------------------------------------------------------------------------
%	CONTENT
%-------------------------------------------------------------------------------
\begin{cventries}
%---------------------------------------------------------
\cventry
{Group Project} % Affiliation/role
{\href{https://github.com/joel98ferreira/pi-project-mobile-app}{ExpiredUM \faLink}} % Organization/group
{Braga, Portugal} % Location
{2021 - 2022} % Date(s)
{
  \begin{cvitems} % Description(s) of experience/contributions/knowledge
    \item {University group project to create a Mobile app to keep cataloged everything in our home.}
    \item {React Native - Java - Spring Boot such as ASP.NET, C\# and MySQL for the database, and {Microsoft Azure tools} for cloud}
  \end{cvitems}
}
%---------------------------------------------------------
\cventry
{Individual Project} % Affiliation/role
{\href{https://github.com/filipeguimaraes/VR_TP2}{Simpified OAuth \faLink}} % Organization/group
{Braga, Portugal} % Location
{2021} % Date(s)
{
  \begin{cvitems} % Description(s) of experience/contributions/knowledge
    \item {A authentication server with different roles implemented in containers.}
    \item {Docker - Docker Compose - Node.js - Express.js - MySQL} 
  \end{cvitems}
}
%---------------------------------------------------------
\cventry
{Individual Project} % Affiliation/role
{\href{https://github.com/filipeguimaraes/GR_TP2}{SNMP Monitoring \faLink}} % Organization/group
{Braga, Portugal} % Location
{2021} % Date(s)
{
  \begin{cvitems} % Description(s) of experience/contributions/knowledge
    \item {Hosts hardware monitoring using SNMP.}
    \item {NET-SNMP - INMF - MIBs - JAVA - SNMP4j - JavaFx} 
  \end{cvitems}
}
%---------------------------------------------------------
    \cventry
      {Group Project} % Affiliation/role
      {\href{https://github.com/filipeguimaraes/LI4_1920}{SportsManager \faLink}} % Organization/group
      {Braga, Portugal} % Location
      {2020} % Date(s)
      {
        \begin{cvitems} % Description(s) of experience/contributions/knowledge
          \item {University group project to create a web app where gym lessons and courts are available for renting.}
          \item {Front-end ReactJs - ASP.NET - C\# - MySQL - Microsoft Azure tools for cloud}
        \end{cvitems}
      }
%---------------------------------------------------------
 
%---------------------------------------------------------
\end{cventries}

\end{minipage}
%-------------------------------------------------------------------------------
%	SECTION TITLE
%-------------------------------------------------------------------------------
\cvsection{Skills}


%-------------------------------------------------------------------------------
%	CONTENT
%-------------------------------------------------------------------------------
  
\begin{cvskills}

%---------------------------------------------------------
  \cvskill
    {Good} % Category
    {React, React Native, Expo, HTML5, JavaFx, Express/Node, Java, JAvaScript, Docker, NGINX, HLS, LATEX }
%---------------------------------------------------------
  \cvskill
    {Intermediary} % Category
    {Kubernetes, CSS, Node.js, C, C++, OpenGL, MySQL} % Skills

%---------------------------------------------------------
  \cvskill
    {Basic} % Category
    {OpenFlow, Ansible, Spring Boot, Haskell, C#}

  \cvskill
    {Languages} % Category
    {Portuguese, English} % Skills

%---------------------------------------------------------
\end{cvskills}

%%-------------------------------------------------------------------------------
%	SECTION TITLE
%-------------------------------------------------------------------------------
\cvsection{Honors \& Awards}


%-------------------------------------------------------------------------------
%	SUBSECTION TITLE
%-------------------------------------------------------------------------------
\cvsubsection{International}


%-------------------------------------------------------------------------------
%	CONTENT
%-------------------------------------------------------------------------------
\begin{cvhonors}

%---------------------------------------------------------
  \cvhonor
    {Finalist} % Award
    {DEFCON 26th CTF Hacking Competition World Final} % Event
    {Las Vegas, U.S.A} % Location
    {2018} % Date(s)

%---------------------------------------------------------
  \cvhonor
    {Finalist} % Award
    {DEFCON 25th CTF Hacking Competition World Final} % Event
    {Las Vegas, U.S.A} % Location
    {2017} % Date(s)

%---------------------------------------------------------
  \cvhonor
    {Finalist} % Award
    {DEFCON 22nd CTF Hacking Competition World Final} % Event
    {Las Vegas, U.S.A} % Location
    {2014} % Date(s)

%---------------------------------------------------------
  \cvhonor
    {Finalist} % Award
    {DEFCON 21st CTF Hacking Competition World Final} % Event
    {Las Vegas, U.S.A} % Location
    {2013} % Date(s)

%---------------------------------------------------------
  \cvhonor
    {Finalist} % Award
    {DEFCON 19th CTF Hacking Competition World Final} % Event
    {Las Vegas, U.S.A} % Location
    {2011} % Date(s)

%---------------------------------------------------------
  \cvhonor
    {6th Place} % Award
    {SECUINSIDE Hacking Competition World Final} % Event
    {Seoul, S.Korea} % Location
    {2012} % Date(s)

%---------------------------------------------------------
\end{cvhonors}


%-------------------------------------------------------------------------------
%	SUBSECTION TITLE
%-------------------------------------------------------------------------------
\cvsubsection{Domestic}


%-------------------------------------------------------------------------------
%	CONTENT
%-------------------------------------------------------------------------------
\begin{cvhonors}

%---------------------------------------------------------
  \cvhonor
    {3rd Place} % Award
    {WITHCON Hacking Competition Final} % Event
    {Seoul, S.Korea} % Location
    {2015} % Date(s)

%---------------------------------------------------------
  \cvhonor
    {Silver Prize} % Award
    {KISA HDCON Hacking Competition Final} % Event
    {Seoul, S.Korea} % Location
    {2017} % Date(s)

%---------------------------------------------------------
  \cvhonor
    {Silver Prize} % Award
    {KISA HDCON Hacking Competition Final} % Event
    {Seoul, S.Korea} % Location
    {2013} % Date(s)

%---------------------------------------------------------
  \cvhonor
    {2nd Award} % Award
    {HUST Hacking Festival} % Event
    {S.Korea} % Location
    {2013} % Date(s)

%---------------------------------------------------------
  \cvhonor
    {3rd Award} % Award
    {HUST Hacking Festival} % Event
    {S.Korea} % Location
    {2010} % Date(s)

%---------------------------------------------------------
  \cvhonor
    {3rd Award} % Award
    {Holyshield 3rd Hacking Festival} % Event
    {S.Korea} % Location
    {2012} % Date(s)

%---------------------------------------------------------
  \cvhonor
    {2nd Award} % Award
    {Holyshield 3rd Hacking Festival} % Event
    {S.Korea} % Location
    {2011} % Date(s)

%---------------------------------------------------------
  \cvhonor
    {5th Place} % Award
    {PADOCON Hacking Competition Final} % Event
    {Seoul, S.Korea} % Location
    {2011} % Date(s)

%---------------------------------------------------------
\end{cvhonors}

%%-------------------------------------------------------------------------------
%	SECTION TITLE
%-------------------------------------------------------------------------------
\cvsection{Presentation}


%-------------------------------------------------------------------------------
%	CONTENT
%-------------------------------------------------------------------------------
\begin{cventries}

%---------------------------------------------------------
  \cventry
    {Presenter for <Hosting Web Application for Free utilizing GitHub, Netlify and CloudFlare>} % Role
    {DevFest Seoul by Google Developer Group Korea} % Event
    {Seoul, S.Korea} % Location
    {Nov. 2017} % Date(s)
    {
      \begin{cvitems} % Description(s)
        \item {Introduced the history of web technology and the JAM stack which is for the modern web application development.}
        \item {Introduced how to freely host the web application with high performance utilizing global CDN services.}
      \end{cvitems}
    }

%---------------------------------------------------------
  \cventry
    {Presenter for <DEFCON 20th : The way to go to Las Vegas>} % Role
    {6th CodeEngn (Reverse Engineering Conference)} % Event
    {Seoul, S.Korea} % Location
    {Jul. 2012} % Date(s)
    {
      \begin{cvitems} % Description(s)
        \item {Introduced CTF(Capture the Flag) hacking competition and advanced techniques and strategy for CTF}
      \end{cvitems}
    }

%---------------------------------------------------------
  \cventry
    {Presenter for <Metasploit 101>} % Role
    {6th Hacking Camp - S.Korea} % Event
    {S.Korea} % Location
    {Sep. 2012} % Date(s)
    {
      \begin{cvitems} % Description(s)
        \item {Introduced basic procedure for penetration testing and how to use Metasploit}
      \end{cvitems}
    }

%---------------------------------------------------------
\end{cventries}

%%-------------------------------------------------------------------------------
%	SECTION TITLE
%-------------------------------------------------------------------------------
\cvsection{Writing}


%-------------------------------------------------------------------------------
%	CONTENT
%-------------------------------------------------------------------------------
\begin{cventries}

%---------------------------------------------------------
  \cventry
    {Founder \& Writer} % Role
    {A Guide for Developers in Start-up} % Title
    {Facebook Page} % Location
    {Jan. 2015 - PRESENT} % Date(s)
    {
      \begin{cvitems} % Description(s)
        \item {Drafted daily news for developers in Korea about IT technologies, issues about start-up.}
      \end{cvitems}
    }

%---------------------------------------------------------
  \cventry
    {Undergraduate Student Reporter} % Role
    {AhnLab} % Title
    {S.Korea} % Location
    {Oct. 2012 - Jul. 2013} % Date(s)
    {
      \begin{cvitems} % Description(s)
        \item {Drafted reports about IT trends and Security issues on AhnLab Company magazine.}
      \end{cvitems}
    }

%---------------------------------------------------------
\end{cventries}

%%-------------------------------------------------------------------------------
%	SECTION TITLE
%-------------------------------------------------------------------------------
\cvsection{Program Committees}


%-------------------------------------------------------------------------------
%	CONTENT
%-------------------------------------------------------------------------------
\begin{cvhonors}

%---------------------------------------------------------
  \cvhonor
    {Problem Writer} % Position
    {2016 CODEGATE Hacking Competition World Final} % Committee
    {S.Korea} % Location
    {2016} % Date(s)

%---------------------------------------------------------
  \cvhonor
    {Organizer \& Co-director} % Position
    {1st POSTECH Hackathon} % Committee
    {S.Korea} % Location
    {2013} % Date(s)

%---------------------------------------------------------
  \cvhonor
    {Staff} % Position
    {7th Hacking Camp} % Committee
    {S.Korea} % Location
    {2012} % Date(s)

%---------------------------------------------------------
  \cvhonor
    {Problem Writer} % Position
    {1st Hoseo University Teenager Hacking Competition} % Committee
    {S.Korea} % Location
    {2012} % Date(s)

%---------------------------------------------------------
  \cvhonor
    {Staff \& Problem Writer} % Position
    {JFF(Just for Fun) Hacking Competition} % Committee
    {S.Korea} % Location
    {2012} % Date(s)

%---------------------------------------------------------
\end{cvhonors}



%-------------------------------------------------------------------------------
\end{document}
